\textbf{1. Kiến thức cần nhớ}
\begin{itemize}
\item \textbf{So sánh độ dài của đường kính và dây}: Trong các dây của một đường tròn, dây lớn nhất là \textbf{đường kính}:
$$\boxed{AB \le 2R}$$
\item \textbf{Quan hệ vuông góc giữa đường kính và dây}:
\begin{itemize}
\item Đường kính vuông góc với một dây thì \textbf{đi qua trung điểm} của dây ấy.
\item Đường kính đi qua trung điểm của một dây \textbf{không đi qua tâm} thì \textbf{vuông góc} với dây ấy.
\end{itemize}
\item \textbf{Liên hệ giữa dây và khoảng cách từ tâm đến dây}: Trong một đường tròn:
\begin{itemize}
\item Hai dây bằng nhau thì \textbf{cách đều tâm} ($AB = CD \iff OH = OK$).
\item Hai dây cách đều tâm thì \textbf{bằng nhau}.
\item Dây nào lớn hơn thì dây đó \textbf{gần tâm hơn} ($AB > CD \iff OH < OK$).
\end{itemize}
\item \textbf{Góc ở tâm}: Là góc có đỉnh trùng với tâm đường tròn.
$$\boxed{\text{Số đo cung nhỏ } AB = \text{Số đo góc ở tâm } \widehat{AOB}}$$
\end{itemize}

\textbf{2. Phương pháp giải các dạng}

%% DANG 9H5?14-1
\textbf{Dạng 1. Quan hệ giữa đường kính và dây}
\begin{enumerate}
\item Nhận diện tam giác vuông tạo bởi tâm $O$, trung điểm $H$ của dây $AB$ và một đầu mút $A$ của dây:
$$\triangle OHA\text{ vuông tại } H \implies \boxed{R^2 = OH^2 + \left(\dfrac{AB}{2}\right)^2}$$
\item Biết hai trong ba đại lượng ($R$, khoảng cách $OH$, độ dài dây $AB$) tính đại lượng còn lại.
\end{enumerate}

%% DANG 9H5?14-2
\textbf{Dạng 2. Góc ở tâm, số đo cung}
\begin{enumerate}
\item Số đo của cung nhỏ bằng số đo của góc ở tâm chắn cung đó: $\text{sđ}\wideparen{AB} = \widehat{AOB}$.
\item Số đo của cung lớn bằng $360^\circ$ trừ đi số đo của cung nhỏ: $\text{sđ}\wideparen{AnB} = 360^\circ - \text{sđ}\wideparen{AmB}$.
\item Số đo của nửa đường tròn bằng $180^\circ$.
\end{enumerate}

%% DANG 9H5?14-3
\textbf{Dạng 3. Liên hệ giữa dây và khoảng cách đến tâm}
\begin{enumerate}
\item Để so sánh độ dài hai dây: so sánh khoảng cách từ tâm đến hai dây (dây gần tâm hơn thì lớn hơn).
\item Để chứng minh hai dây bằng nhau: chứng minh khoảng cách từ tâm đến hai dây bằng nhau ($OH = OK$).
\end{enumerate}
